% --- Archon physics formalization source begin ---
% archon:physics
% archon:chemistry
% archon:covers IChO2026Problems/problem_icho_2026_t9_a7.lean
% archon:source-report reports/icho_2026/problem_icho_2026_t9_a7.source.json
% archon:problem-id icho_2026_t9
% archon:part-id T9-A7
% archon:previous-part icho_2026_t9_a5
% archon:previous-part-policy natural_language_prerequisite_only

\section{Icho Chemistry problem icho\_2026\_t9\_a7}

\paragraph{Problem source.}
\#\# Source
58th International Chemistry Olympiad, Tashkent, Uzbekistan, 2026.
Official-materials index: https://www.icho2026.uz/problems
English problem PDF: ../icho\_2026\_source/raw/theory\_problem.pdf
English solution/rubric PDF: ../icho\_2026\_source/raw/theory\_solution.pdf
Problem PDF page 87; printed local question page 4; solution starts on PDF page 87.
The page PNG is primary visual evidence for formulas, structures, tables, and figures.

\#\# T9. Cyclodextrin Chemistry
T9. Cyclodextrin Chemistry 7\% Cyclodextrins (CD) are a family of cyclic oligosaccharides, consisting of glucose subunits joined by α1,4-glycosidic bonds. The three most common cyclodextrins (α-, β-, γ-cyclodextrin) contain 6, 7, and 8 α-D-glucopyranoside units, respectively. 9.1 Calculate the molar mass of β-CD. Assume M𝑤(glucose) = 180.16 g mol−1. CDs combine macrocyclic properties with glucose chemistry, which was demonstrated in the synthesis of X. Primary and secondary hydroxy groups show different reactivity due to being a part of CDs. Stoddart et al. synthesised compound K, where only one OH group remained free per glucopyranoside unit. equiv.= equivalent 9.2 Tick the favorable chair conformation and draw the structure of K by completing the CD template. Show stereochemistry unambiguously. 9.3 Determine the ring size, 𝑟𝑠, of macrocycle X. Give the number of stereocentres, 𝑠𝑐, in X. α-Cycloaltrin, a synthetic analogue of α-CD, was synthesised as follows. 9.4 Draw the structure of Y with stereochemistry by completing the CD template. In 2000, Sinay et al. demonstrated that a single protic group (NH and OH) at unit 1 directs the next reductive debenzylation of a primary OH group to unit 4 in the macrocyclic ring (or to unit 3 if the unit 4 position is not available). This allowed the synthesis of the following β-CD dimer as a mixture of constitutional linkage isomers. 9.5 Draw the structure of L on the β-CD template. Fill in all the boxes. 9.6 Determine the number of isomers of the β-CD dimer that can form during the synthesis by Sinay et. al. To confirm the positions of debenzylation, L was subjected to the “hexo-5-enose degradation” reaction, as shown below. The products were then analysed by mass spectrometry.

\#\# Current subquestion 9.7
Calculate the 𝑚/𝑧value for the two [M+Na]+ peaks that were observed for the degradation products from L. Use integer values of atomic mass.

\#\# Dataset metadata
IChO 2026; T9; raw rubric points=6; kind=theory; formalization\_ready=true.

\paragraph{Shared problem context.}
T9. Cyclodextrin Chemistry 7\% Cyclodextrins (CD) are a family of cyclic oligosaccharides, consisting of glucose subunits joined by α1,4-glycosidic bonds. The three most common cyclodextrins (α-, β-, γ-cyclodextrin) contain 6, 7, and 8 α-D-glucopyranoside units, respectively. 9.1 Calculate the molar mass of β-CD. Assume M𝑤(glucose) = 180.16 g mol−1. CDs combine macrocyclic properties with glucose chemistry, which was demonstrated in the synthesis of X. Primary and secondary hydroxy groups show different reactivity due to being a part of CDs. Stoddart et al. synthesised compound K, where only one OH group remained free per glucopyranoside unit. equiv.= equivalent 9.2 Tick the favorable chair conformation and draw the structure of K by completing the CD template. Show stereochemistry unambiguously. 9.3 Determine the ring size, 𝑟𝑠, of macrocycle X. Give the number of stereocentres, 𝑠𝑐, in X. α-Cycloaltrin, a synthetic analogue of α-CD, was synthesised as follows. 9.4 Draw the structure of Y with stereochemistry by completing the CD template. In 2000, Sinay et al. demonstrated that a single protic group (NH and OH) at unit 1 directs the next reductive debenzylation of a primary OH group to unit 4 in the macrocyclic ring (or to unit 3 if the unit 4 position is not available). This allowed the synthesis of the following β-CD dimer as a mixture of constitutional linkage isomers. 9.5 Draw the structure of L on the β-CD template. Fill in all the boxes. 9.6 Determine the number of isomers of the β-CD dimer that can form during the synthesis by Sinay et. al. To confirm the positions of debenzylation, L was subjected to the “hexo-5-enose degradation” reaction, as shown below. The products were then analysed by mass spectrometry.

\paragraph{Current subquestion.}
Calculate the 𝑚/𝑧value for the two [M+Na]+ peaks that were observed for the degradation products from L. Use integer values of atomic mass.

\paragraph{Recorded answer/context.}
Since we know the chemical formula for the units, we can calculate molecular weights of the two products. Each product should have one unit 2 (black sphere) and 2-3 unit 1s (white sphere). Additionally, we should add C2H3O2, which corresponds to the mass of the OAc group. Therefore: 𝑀1 ∶2C27H28O5 + C22H25O4 + C2H3O2 = C78H84O16 [𝑀1 + Na]+= 1299 𝑀2 ∶3C27H28O5 + C22H25O4 + C2H3O2 = C105H112O21 [𝑀2 + Na]+=1731 [𝑀1 + Na]+ = 1299 [𝑀2 + Na]+= 1731 Each correct answer 3 points, total – 6 points m/z without Na+ 1p each m/z without C2H2O2 1p each 6.0 pt

\paragraph{Figure/image paths.}
\begin{itemize}
\item ../icho\_2026\_source/image/T9\_page-4.png
\item ../icho\_2026\_source/image/T9\_page-3.png
\end{itemize}

\paragraph{Reusable previous-part conclusions.}
\begin{itemize}
\item Source T9-A5. Question: Draw the structure of L on the β-CD template. Fill in all the boxes. Reusable conclusions: Visual-answer notice: the official solution contains essential ticks, structures, tables, graphs, or equations that PDF plain-text extraction may omit. Consult the cited solution PDF page.

Correct structure 2 points. -1 p for each incorrect substituent 0 p if 2 or more incorrect substituents 2.0 pt Policy: natural\_language\_prerequisite\_only; the current target and later answers must not be assumed
\end{itemize}

\paragraph{Formalization target.}
create a compiling Lean file with sorry bodies at `IChO2026Problems/problem\_icho\_2026\_t9\_a7.lean`.
The Lean declarations must preserve every mathematical object, quantifier, hypothesis, side condition, approximation/error guarantee, and requested conclusion from the source statement.
Use Mathlib/Physlib/Chemistry names found through LeanExplore where available. Prefer the configured benchmark Base library; introduce a local definition only when the source genuinely requires an object that the environment does not expose.

\begin{theorem}\leanok
[Icho Chemistry formalization target]
\label{thm:physics:icho_2026_t9_a7:target}
\lean{IChO2026Problems.T9A7.sodium_adduct_peaks_of_L_degradation_products}
The degradation-product formulas and sodium-adduct convention determine the
mass-spectrometric peaks of L.
\end{theorem}
\begin{proof}

\leanok
The two source formulas are respectively $\mathrm{C}_{78}\mathrm{H}_{84}
\mathrm{O}_{16}$ and $\mathrm{C}_{105}\mathrm{H}_{112}\mathrm{O}_{21}$.
Using the required integer masses gives neutral nominal masses
\(78\cdot12+84+16\cdot16=1276\) and
\(105\cdot12+112+21\cdot16=1708\).  A singly charged sodium adduct adds one
sodium mass, $23$, and division by charge one changes neither total.  Thus
the two observed $m/z$ values are $1276+23=1299$ and $1708+23=1731$.
\end{proof}
% --- Archon physics formalization source end ---
